\emph{On note $I$ l'évènement \og le client a acheté sur internet \fg{} et $S$ l'évènement \og le client est satisfait \fg.}

\begin{enumerate}
	\item Acheteurs en magasin : $\np{1000} - 400 = 600$.

	Internet et satisfaits : $0,9 \times 400 = 360$, donc internet et non satisfaits : $400 - 360 = 40$.

	Magasin et non satisfaits : $0,2 \times 600 = 120$, donc magasin et satisfaits : $600 - 120 = 480$.

	\begin{center}
		\begin{tblr}{hlines, vlines, colspec={X[l,4cm]X[c,3.5cm]X[c,3.5cm]X[c,2cm]}}
			& Achat sur internet ($I$) & Achat en magasin ($\overline{I}$) & Total \\
			Satisfait ($S$) & 360 & 480 & 840 \\
			Non satisfait ($\overline{S}$) & 40 & 120 & 160 \\
			Total & 400 & 600 & \np{1000} \\
		\end{tblr}
	\end{center}

	\item Le client étant choisi au hasard, on est en situation d'équiprobabilité. La probabilité que le client ait acheté en magasin et soit satisfait est $P\left(\overline{I} \cap S\right)$.
	
	\[P\left(\overline{I} \cap S\right) = \dfrac{480}{\np{1000}} = 0,48.\]

La probabilité que le client ait acheté en magasin et soit satisfait est égale à $0,48$.

	\item Parmi les 160 clients non satisfaits, 40 ont acheté sur internet :
	\[P_{\overline{S}}(I) = \dfrac{40}{160} = 0,25.\]
	Sachant que le client n'est pas satisfait, la probabilité qu'il ait acheté sur internet est égale à
	$0,25$.

	\item La proportion de clients non satisfaits est
	$\dfrac{160}{\np{1000}} = 0,16 = 16\,\%$.

	Or $16\,\% > 10\,\%$ : la stratégie d'amélioration sera mise en place.
\end{enumerate}
